\newpage    
\thispagestyle{empty}
\phantomsection
\addtotoc{Abstract}

\begin{center}\huge \textbf{Abstract}\end{center}

This thesis presents the full-custom design and implementation of a PCI Express (PCIe)
Gen~6.0 Serializer/Deserializer (SerDes) Physical Layer (PHY) Transceiver operating at
64~GT/s in GlobalFoundries 22\,nm FD-SOI (GF22FDX) technology. The transition from
NRZ to PAM4 signaling introduces profound signal integrity challenges, including
$\sim$26\,dB channel insertion loss at Nyquist, severe inter-symbol interference (ISI),
and stringent linearity requirements across all PVT corners.

On the transmit side, a quadrature clocking architecture with IQ frequency divider and
25\% duty-cycle pulse-shaper minimizes clock distribution jitter. A 4:1 C\textsuperscript{2}MOS
multiplexer achieves a post-layout eye width of 30.35\,ps under worst-case (SF/125\textdegree{}C)
conditions and 155.29\,fs deterministic jitter at the typical corner, well within the
625\,fs PCIe Gen~6.0 budget. The line driver employs a segmented 32\,GBaud Voltage-Mode
SST DAC with peaking inductors, delivering $>$800\,mV differential swing at 10.22\,mW
(PAM4, typical corner), with a multi-tap FFE cancelling precursor ISI.

On the receive side, T-coil and bridged-shunt peaking networks extend AFE bandwidth
while preserving impedance matching. A source-degenerated VGA maximizes dynamic range,
followed by a multi-stage CTLE realizing a 6-pole, 3-zero response with up to 19\,dB
programmable peaking to counteract channel loss. A half-rate look-ahead DFE---comprising
an Azita strong-ARM comparator, C\textsuperscript{2}MOS latches, G\textsubscript{m}
cells, and an analog dLev calibration loop---cancels post-cursor ISI, achieving
60--66\,mV eye height and 2.88\,mV\textsubscript{rms} integrated noise over a
26.5\,dB FR4 channel. The complete PHY is validated across 45 PVT corners, confirming
compliance with PCIe Gen~6.0 electrical specifications.